
\documentclass{article}
\usepackage{colortbl}
\usepackage{makecell}
\usepackage{multirow}
\usepackage{supertabular}

\begin{document}

\newcounter{utterance}

\centering \large Interaction Transcript for game `ifeval', experiment `ifeval', episode 37 with SLED{-}BSU\_\_Qwen3.5{-}27B{-}sft{-}ep1\_\_prm{-}ep1{-}best\_of\_n.
\vspace{24pt}

{ \footnotesize  \setcounter{utterance}{1}
\setlength{\tabcolsep}{0pt}
\begin{supertabular}{c@{$\;$}|p{.15\linewidth}@{}p{.15\linewidth}p{.15\linewidth}p{.15\linewidth}p{.15\linewidth}p{.15\linewidth}}
    \# & \multicolumn{2}{c}{Player} && \multicolumn{2}{c}{Game Master} \\
    \hline

    \theutterance \stepcounter{utterance}  
    & & & \multicolumn{4}{p{0.6\linewidth}}{
        \cellcolor[rgb]{0.9,0.9,0.9}{
            \makecell[{{p{\linewidth}}}]{
                \texttt{\tiny{[P1$\langle$GM]}}
                \texttt{What sentiments existed in the Croatian and Serbian nations towards the Austrian Empire at the end of the 19th century? Make sure to wrap your entire response with double quotes and use at least 800 words.} \\
            }
        }
    }
     \\ \\

    \theutterance \stepcounter{utterance}  
    & \multicolumn{4}{p{0.6\linewidth}}{
        \cellcolor[rgb]{0.561,0.737,0.561}{
            \makecell[{{p{\linewidth}}}]{
                \texttt{\tiny{[P1$\rangle$GM]}}
                \texttt{"By the end of the nineteenth century, the sentiments of the Croatian and Serbian nations toward the Austrian Empire were complex, contradictory, and deeply rooted in the divergent political, cultural, and religious realities of the two groups. While both nations shared a common language and a long history of Slavic solidarity, their differing geopolitical positions and aspirations led to markedly different attitudes toward Vienna’s rule. For Croats, the relationship was one of cautious cooperation tempered by frustration, whereas for Serbs, it was largely one of resistance and suspicion, shaped by their experience under direct Ottoman and later Austro{-}Hungarian domination.} \\
\\ 
\texttt{For the Croatian nation, the Austrian Empire—and specifically the Habsburg Monarchy—represented a double{-}edged sword. On one hand, Croatia had been part of the Habsburg realm since the early modern period, and by the late nineteenth century, it enjoyed a degree of autonomy within the Austro{-}Hungarian framework. The Croatian{-}Slavonian Banate, established after the Austro{-}Hungarian Compromise of 1867, granted Croatia internal self{-}government in matters of education, justice, and local administration. This autonomy was seen by many Croatian nationalists as a hard{-}won achievement that protected their language, culture, and political identity against both Hungarian and German dominance. Consequently, a significant segment of the Croatian elite viewed Vienna as a necessary counterweight to Budapest’s attempts at Magyarization. The Croatian Party of Rights, for instance, while critical of Hungarian encroachment, often sought to strengthen ties with Vienna to secure greater autonomy or even a tripartite structure within the empire that would elevate Croatia to equal status with Austria and Hungary.} \\
\\ 
\texttt{Moreover, the Croatian national revival, or Illyrian movement, had flourished under Habsburg rule, benefiting from the relative stability and institutional support that the empire provided. The use of the Croatian language in schools, the press, and official documents was permitted and even encouraged by Vienna as a means to dilute Hungarian influence. This fostered a sense of loyalty among many Croats, who saw the Habsburgs as patrons of Slavic culture and defenders of their national rights. The Croatian national anthem, 'Lijepa naša domovino,' composed in the 1840s, was a testament to this patriotic sentiment, which found expression in the broader framework of the empire.} \\
\\ 
\texttt{However, this loyalty was not unconditional. Many Croats were acutely aware of their subordinate position within the empire and feared that their autonomy was fragile and subject to the whims of Hungarian politics. The failure of the Croatian{-}Hungarian agreement of 1868 to grant Croatia full parity with Hungary, and the subsequent erosion of their rights in the late nineteenth century, led to growing disillusionment. The Croatian national movement increasingly turned toward the idea of a South Slav federation, which would unite Croats, Serbs, and Slovenes in a single state, either as a third partner in the empire or as an independent nation. This shift marked a gradual move away from loyalty to Vienna toward a more assertive national agenda.} \\
\\ 
\texttt{In contrast, the Serbian nation’s relationship with the Austrian Empire was far more fraught. Unlike the Croats, the Serbs did not enjoy autonomy within the Habsburg realm. The majority of Serbs lived in the Kingdom of Serbia, which had achieved formal independence from the Ottoman Empire in 1878, or in the Principality of Montenegro, both of which were sovereign states outside the empire. For these Serbs, the Austrian Empire was a foreign power, often viewed with suspicion as a potential threat to their hard{-}won independence. The Serbs had a long history of resistance to Habsburg rule, dating back to the Great Migrations of the seventeenth century, and by the late nineteenth century, they were keenly aware of the empire’s expansionist ambitions in the Balkans.} \\
\\ 
\texttt{For the smaller number of Serbs who lived within the Austrian Empire, particularly in the Military Frontier and the crown land of Bosnia{-}Herzegovina (occupied by Austria{-}Hungary in 1878), the situation was even more precarious. In Bosnia, the Serbs were a religious and cultural minority under Austro{-}Hungarian administration, which they perceived as hostile to their Orthodox identity and national aspirations. The Habsburgs’ preference for Catholicism and their promotion of a separate Bosnian national identity were seen as deliberate attempts to undermine Serbian unity and loyalty. This led to widespread resentment among Serbs in Bosnia, who often aligned themselves with the Kingdom of Serbia and advocated for unification with their co{-}nationals outside the empire.} \\
\\ 
\texttt{The assassination of Archduke Franz Ferdinand in 1914, which triggered World War I, was a direct result of these tensions. The assassin, Gavrilo Princip, was a Bosnian Serb nationalist who sought to liberate his people from Austro{-}Hungarian rule and unite them with Serbia. This act was not an isolated incident but the culmination of decades of Serbian nationalist agitation against the empire. For the Serbian nation as a whole, the Austrian Empire was seen as an obstacle to the realization of their national goals, whether that meant independence, expansion, or the creation of a Greater Serbia.} \\
\\ 
\texttt{In summary, the sentiments of Croats and Serbs toward the Austrian Empire at the end of the nineteenth century were shaped by their distinct historical experiences and political positions. For Croats, the empire was a mixed blessing—a protector of their autonomy and culture but also a source of frustration and limitation. For Serbs, it was largely a source of oppression and a barrier to their national aspirations. These divergent attitudes would play a crucial role in the unfolding of events in the early twentieth century, ultimately contributing to the collapse of the Habsburg Monarchy and the redrawing of the map of the Balkans."} \\
            }
        }
    }
    & & \\ \\

    \theutterance \stepcounter{utterance}  
    & & & \multicolumn{2}{p{0.3\linewidth}}{
        \cellcolor[rgb]{0.9,0.9,0.9}{
            \makecell[{{p{\linewidth}}}]{
                \texttt{\tiny{[GM$|$GM]}}
                {'startend:quotation': {}, 'length_constraints:number_words': {'relation': 'at least', 'num_words': 800}}
            }
        }
    }
    & & \\ \\

    \theutterance \stepcounter{utterance}  
    & & & \multicolumn{2}{p{0.3\linewidth}}{
        \cellcolor[rgb]{0.9,0.9,0.9}{
            \makecell[{{p{\linewidth}}}]{
                \texttt{\tiny{[GM$|$GM]}}
                \texttt{game\_result = WIN} \\
            }
        }
    }
    & & \\ \\

\end{supertabular}
}

\end{document}
